\documentclass[11pt]{amsart}
\usepackage{amssymb}
\usepackage[margin=1.15in]{geometry}
\usepackage[hidelinks]{hyperref}
\hypersetup{pdftitle={Knus-Ojanguren 3.2 for twisted modules}}

\newtheorem{theorem}{Theorem}[section]
\newtheorem{lemma}[theorem]{Lemma}
\newtheorem{proposition}[theorem]{Proposition}
\newtheorem{corollary}[theorem]{Corollary}
\theoremstyle{definition}
\newtheorem{definition}[theorem]{Definition}
\theoremstyle{remark}
\newtheorem{remark}[theorem]{Remark}
\newcommand{\Spec}{\operatorname{Spec}}
\newcommand{\Hom}{\operatorname{Hom}}
\newcommand{\End}{\operatorname{End}}
\newcommand{\GL}{\operatorname{GL}}
\newcommand{\rk}{\operatorname{rk}}
\newcommand{\Q}{\mathbb Q}
\newcommand{\Kt}{\widetilde K_0}

\title{Knus--Ojanguren's recollement (3.2) for twisted modules}
\date{November 2, 2026}

% stamp on every page: file name and date/time of compilation
\newcount\stamphour \newcount\stampmin
\stamphour=\time \divide\stamphour by 60
\stampmin=\time \multiply\stamphour by 60 \advance\stampmin by -\stamphour \divide\stamphour by 60
\edef\stamptime{\the\year-\ifnum\month<10 0\fi\the\month-\ifnum\day<10 0\fi\the\day\space
 \ifnum\stamphour<10 0\fi\the\stamphour:\ifnum\stampmin<10 0\fi\the\stampmin}
\AddToHook{shipout/foreground}{\put(0.5\paperwidth,-\dimexpr\paperheight-0.45in\relax)%
 {\makebox[0pt]{\scriptsize\ttfamily\jobname.pdf\quad\stamptime}}}
\begin{document}
\begin{abstract}
We rewrite Lemma~3.2 of Knus--Ojanguren (LNM 844, p.~224) for twisted
modules and list every ingredient it uses. One ingredient is not stated by
the authors: that some multiple of a class of positive rank in $K_0$ is the
class of a projective module. This fact is not elementary. For the use we
have in mind (gluing over a Zariski cover on whose pieces the gerbe is
trivial) we give a different argument that avoids it. The argument needs
only Lombardi--Quitt\'e X.5.5 ($\Kt$ is nil) and XVI.3.4 (Quillen's
patching lemma), together with elementary manipulations of $K_0$.
\end{abstract}
\maketitle

\section{Setting}

Let $A$ be a commutative ring and $G$ a $\mathbb G_m$-gerbe over $\Spec A$.
A \emph{twisted module} on an open $U$ is a family $V$ of finitely generated
projective modules over $G|_U$ on which each loop $a\in R^\times$ acts as
multiplication by $a$ (weight $1$). Weight-$k$ modules are defined in the
same way. A weight-$0$ family descends to an ordinary module on $U$. In
particular, for twisted $V,W$ the modules $\Hom(V,W)$ and $\End V$ are
ordinary. Twisted modules are glued along $D(a)\cup D(b)$ like any other
family, since $D(a)\vee D(b)$ is the pushout of $D(a)\leftarrow D(ab)\to
D(b)$.

\emph{Torsion.} The class of $G$ has order dividing $n$ if and only if there
is a weight-$n$ twisted module $\Lambda$ of rank one on all of $\Spec A$. We
fix $\Lambda$, and for a twisted $V$ we set
\[ P(V)=V^{\otimes n}\otimes\Lambda^{-1},\]
an ordinary module of rank $(\rk V)^n$. One has $P(V\otimes F)=P(V)\otimes
F^{\otimes n}$ for $F$ ordinary.

\emph{$K_0$.} We write $K=K_0(U)\otimes\Q$. For $F$ of constant rank $N\geq
1$ put $\rho(F)=[F]/N\in K$. A twisted module $V$ of constant positive rank
is \emph{good} if $\rho(P(V))=1$ in $K$, i.e.\ some $P(V)^{k}$ is stably
free.

\section{KO 3.2 rewritten}

KO 3.2 states: if $c\in H^2(A)$ is torsion and $c|_{A_f}$, $c|_{A_g}$ are
represented by Azumaya algebras, then so is $c$. In twisted form: given
twisted $V$ on $A_f$ and $W$ on $A_g$, with ranks $r$ and $s$, produce a
twisted module on $A$.

\smallskip\noindent\emph{The datum.} On $A_{fg}$ put $H=\Hom(V,W)$ (rank
$rs$) and $E=\End V$ (rank $r^2$). The canonical isomorphism
\begin{equation}\label{datum}
V\otimes H\;\cong\; V\otimes V^\vee\otimes W\;\cong\; W\otimes E
\end{equation}
is KO's datum $(A,B,P,Q,\alpha)$ with $P=H$ and $Q=E$. KO first twist an
arbitrary datum by a line bundle $I\in\operatorname{Pic}(A_{fg})$ so that its
class is $c$, and then use the injectivity of
$\operatorname{Br}(A_f,A_g,A_{fg})\to H^2(A)$ (their Theorem~2.11) to see that
the $n$-th power of the datum is trivial. Here (\ref{datum}) is already a
datum for $G$ itself, so neither step is needed. Its $n$-th power is the
isomorphism
\begin{equation}\label{w}
P(V)\otimes H^{\otimes n}\;\cong\;P(W)\otimes E^{\otimes n},
\end{equation}
which is KO's $w: M\otimes P^{(n)}\cong N\otimes Q^{(n)}$ with $M=P(V)$ and
$N=P(W)$.

\smallskip\noindent\emph{KO's computation.} Let $z\in K$ be the $n$-th root
of $\rho(P(V))^{-1}\rho(H)^{-n}$. Then $z\rho(H)$ is the $n$-th root of
$\rho(P(V))^{-1}$, which comes from $A_f$, and $z\rho(E)$ is the $n$-th root
of $\rho(P(W))^{-1}$, which comes from $A_g$. KO then pick an integer $h$ and
modules $X$ on $A_{fg}$, $P'$ on $A_f$, $Q'$ on $A_g$ with $hz=[X]$,
$hz[H]=[P']$, $hz[E]=[Q']$. Applying Bass's cancellation theorem to
$V\otimes X\otimes H\cong W\otimes X\otimes E$ gives an isomorphism $V\otimes
P'^m\cong W\otimes Q'^m$ for large $m$, and gluing gives the result.

\smallskip\noindent\emph{Ingredients.}
\begin{enumerate}
\item $n$-th roots in $1+\Kt\otimes\Q$, existence and uniqueness. Uniqueness
  is used, without being stated, to identify $z\rho(H)$ with the restriction
  of a root computed on $A_f$. Both follow from LQ X.5.5 ($\Kt$ is the
  nilradical of $K_0$), with the root given by a finite binomial series.
\item \emph{Effectiveness:} for $x\in K_0\otimes\Q$ of positive rank, $hx$ is
  the class of a projective module for suitable $h$. KO quote Bass for this
  in one sentence. It is not elementary. If $L$ is a line bundle with
  $\ell=[L]-1$, $\ell^3=0$ and $1,\ell,\ell^2$ linearly independent in
  $K_0\otimes\Q$, then $8[L]^{1/2}=3+6[L]-[L^2]$,
  and no multiple of it is a combination with nonnegative coefficients of
  powers $[L^k]$, $k\in\mathbb Z$: such a combination would be an
  $\mathbb Z$-valued distribution with mean $1/2$ and variance $0$. So a
  non-split bundle has to be constructed. The classical proof reduces to a
  finitely generated $\mathbb Z$-algebra and applies Serre's splitting
  theorem; constructively this is LQ XIV.3.4 together with the
  dimension bounds of LQ XIII and XIV.
\item Cancellation. Only the weak form below is used, and it is elementary
  (Lemma~\ref{efface}). Moreover $X$ may stay virtual, because
  (\ref{datum}) gives $[V][X][H]=[W][X][E]$ for every $X\in K_0(A_{fg})$.
\item Constant ranks (``sans restriction de la g\'en\'eralit\'e''). In the use
  below all ranks are constant by construction.
\end{enumerate}
Of these, only (2) is hard. The rest of the note shows that it is not needed
for the use we have in mind.

\section{Gluing good twisted modules}

\begin{lemma}\label{roots}
Let $F$ have constant rank $N\geq1$. If $\rho(F)^n=1$ in $K$, then
$\rho(F)=1$.
\end{lemma}
\begin{proof}
Put $\mu=[F]-N$. It is nilpotent by LQ X.5.5. From
$k([F]^n-N^n)=0$ we get $k\mu(c+\mu q)=0$ with $c=nN^{n-1}$ and some $q$.
Then $kc^t\mu=0$ once $\mu^t=0$, by expanding $c^t=((c+\mu q)-\mu q)^t$.
\end{proof}

\begin{lemma}[weak cancellation]\label{efface}
Let $X,Y,T$ be twisted modules with $X\oplus T\cong Y\oplus T$, where
$X$ and $Y$ have positive rank. Then $X^{2m}\cong Y^{2m}$ for some $m$.
\end{lemma}
\begin{proof}
First, $T$ is a direct summand of some $X^m$. The evaluation map
$X\otimes\Hom(X,T)\to T$ is surjective, since this can be checked locally,
where $X=L\otimes F$ with $F$ of positive rank. It splits, because
$\Hom(T,X\otimes\Hom(X,T))\to\End T$ is a surjection of ordinary
finitely generated projective modules. Finally $\Hom(X,T)$ is a summand of
some $O^m$. Similarly, $T$ is a summand of $Y^m$.

Write $X^m\cong T\oplus T'$. Then $X^{m+1}\cong Y\oplus T\oplus T'\cong
Y\oplus X^m$, and by induction $X^{2m}\cong Y^m\oplus X^m$. In the same way
$Y^{2m}\cong X^m\oplus Y^m$.
\end{proof}

\begin{proposition}[gluing good modules]\label{cheap}
Let $a,b\in A$ be comaximal. If $V$ on $A_a$ and $W$ on $A_b$ are good, of
ranks $r$ and $s$, then for some $m$ there is an isomorphism $V^{ms}\cong
W^{mr}$ on $A_{ab}$. Gluing gives a twisted module $V'$ on $A$ whose
restrictions to $A_a$ and $A_b$ are good.
\end{proposition}
\begin{proof}
From $H^{\otimes n}\cong P(V)^\vee\otimes P(W)$ we get $\rho(H)^n=1$. The
dual of a stably free module is stably free, so $\rho(P(V)^\vee)=1$.
Lemma~\ref{roots} then gives $\rho(H)=1$, and likewise $\rho(E)=1$. Hence
$k([H]-rs)=0=k([E]-r^2)$ for some $k$. Tensoring the corresponding stable
isomorphisms with $V$ and with $W$ and using (\ref{datum}) gives
$V^{krs}\oplus T\cong W^{kr^2}\oplus T$ with $T=V^j\oplus W^{j'}$.
Lemma~\ref{efface} gives $V^{2m'krs}\cong W^{2m'kr^2}$; take
$m=2m'kr$. Goodness of the restrictions holds
because $P(V^{e})=P(V)\otimes O^{e^n}$.
\end{proof}

The glued $V'$ need not be good on $A$. Its defect lies in the kernel of
restriction, which we now describe.

\smallskip\noindent\emph{Gluing matrices.} For $g\in\GL_M(A_{ab})$ let $P_g$
be the module glued from $A_a^M$ and $A_b^M$ along $g$. Then $P_{\beta
g\alpha}\cong P_g$ for $\beta\in\GL_M(A_b)$ and $\alpha\in\GL_M(A_a)$, and
$P_g\otimes P_h=P_{g\otimes h}$.

\begin{lemma}[homotopy]\label{homotopy}
Let $\varepsilon(T)\in\GL_M(A_{ab}[T])$ with $\varepsilon(0)=I$. Then
$P_{g\varepsilon(1)}\cong P_g$.
\end{lemma}
\begin{proof}
Apply LQ XVI.3.4 to $\varepsilon$ over $(A_a)_b$, and to
$\phi=g\varepsilon g^{-1}$ over $(A_b)_a$. If $U(T,Y)$ works for $s$ in
XVI.3.4, then $U(T,s'Y)$ works for $ss'$, so one $N$ serves both cases. For
large $N$ this gives
$\varepsilon(T+b^NU)\varepsilon(T)^{-1}\in\GL_M(A_a[T,U])$ and
$\phi(T+a^NU)\phi(T)^{-1}\in\GL_M(A_b[T,U])$. Write
$1=a^Nx+b^Ny$. Then
\[ g\varepsilon(1)=\phi(1)g
 =\bigl[\phi(b^Ny+a^Nx)\phi(b^Ny)^{-1}\bigr]\,g\,\varepsilon(b^Ny),\]
where the bracket is the specialisation $T=b^Ny$, $U=x$ of the second
matrix, so it lies in $\GL_M(A_b)$, and $\varepsilon(b^Ny)$ is the
specialisation $T=0$, $U=y$ of the first, so it lies in $\GL_M(A_a)$; this
uses $\varepsilon(0)=I$.
\end{proof}

\begin{lemma}\label{split}
Let $g,h\in\GL_M(A_{ab})$ and $g'\in\GL_{M'}(A_{ab})$.
\begin{enumerate}
\item $[P_g]+[P_h]=[P_{gh}]+M$. Hence $[P_{g^{j}}]=j[P_g]-(j-1)M$ for all
  $j\in\mathbb Z$.
\item $([P_g]-M)([P_{g'}]-M')=0$.
\item If $x\in K_0(A)$ restricts to $0$ in $K_0(A_a)$ and in $K_0(A_b)$, then
  $x=[P_g]-M$ for some $g$.
\end{enumerate}
\end{lemma}
\begin{proof}
(1) $P_g\oplus P_h=P_{\mathrm{diag}(g,h)}$ and $\mathrm{diag}(g,h)=
\mathrm{diag}(gh,1)\,\mathrm{diag}(h^{-1},h)$. Whitehead's identity
$\mathrm{diag}(h^{-1},h)=u(h^{-1})\,l(-h)\,u(h^{-1})\,u(-1)\,l(1)\,u(-1)$,
with $u$ and $l$ the block-unipotent matrices, has an evident homotopy to $I$
obtained by multiplying every block by $T$. Apply Lemma~\ref{homotopy}.
(2) Write $g\otimes g'=(g\otimes1)(1\otimes g')$ and apply (1). Up to
constant permutation matrices, $g\otimes1$ is $\mathrm{diag}(g,\dots,g)$, so
$[P_{g\otimes1}]=M'[P_g]$, and likewise $[P_{1\otimes g'}]=M[P_{g'}]$.
(3) Write $x=[F]-K$. After adding a common free summand, $F\oplus O^j$ is
free of rank $M$ on both pieces. Take $g$ to be the transition between the
two bases.
\end{proof}

\begin{proposition}[correction]\label{correct}
Let $V'$ be a twisted module on $A$ of constant rank whose restrictions to
$A_a$ and $A_b$ are good. Then there is an ordinary $S$ of constant positive
rank such that $V'\otimes S$ is good.
\end{proposition}
\begin{proof}
Let $N=\rk P(V')$ and $\tau=\rho(P(V'))-1$. Since $P$ commutes with
restriction, goodness on $A_a$ and $A_b$ says that $[P(V')]-N$ restricts to
torsion classes there, so $\kappa([P(V')]-N)$ restricts to $0$ on both for
some $\kappa\ge1$. By Lemma~\ref{split}(3), $\kappa N\tau=[P_g]-M$ for some
$g$, and $\tau^2=0$ by
Lemma~\ref{split}(2). Choose $e$ with $R=n\kappa Ne\geq M$ and put
$S=P_{g^{-e}}\oplus O^{R-M}$. By Lemma~\ref{split}(1),
$[S]=R-e([P_g]-M)$, so $\rho(S)=1-\tau/n$. Hence
\[\rho(P(V'\otimes S))=(1+\tau)(1-\tau/n)^n=(1+\tau)(1-\tau)=1.\qedhere\]
\end{proof}

\begin{theorem}\label{main}
Let $u_1,\dots,u_k\in A$ be comaximal, and suppose that on each $D(u_i)$
there is a rank-one twisted $L_i$ with $L_i^{\otimes n}\cong\Lambda$. Then
there is a twisted module of constant positive rank on $\Spec A$.
\end{theorem}
\begin{proof}
Let $\Gamma$ be the set of $s\in A$ such that some good twisted module
exists on $A_s$. It contains each $u_i$, since $P(L_i)\cong O$, and it is
closed under multiplication by elements of $A$. If $s,s'\in\Gamma$, apply
Propositions~\ref{cheap} and~\ref{correct} over $A_{s+s'}$, where $s$ and
$s'$ are comaximal. This gives $s+s'\in\Gamma$. Hence $\Gamma$ is an ideal
and $1\in\Gamma$.
\end{proof}

The hypothesis $L_i^{\otimes n}\cong\Lambda$ costs nothing: the line bundle
$L_i^{\otimes n}\otimes\Lambda^{-1}$ is trivial on a finer Zariski cover.

\section{What this changes for Gabber in Unit}

\begin{enumerate}
\item Effectiveness (ingredient (2)), Bass cancellation and KO's
  Theorem~2.11 are not used. Neither are $n$-th roots beyond the
  square-zero case: only Lemma~\ref{roots} uses LQ X.5.5.
\item Theorem~\ref{main} holds for any ring $A$. In particular it does not
  use that $A$ is semilocal or local-global (LQ IX.6) or the symmetric-group
  trick. The flattening of a $T$-cover already provides a Zariski cover of
  the whole tower $D$ on which $G$ is trivial, so KO's 3.4 step, ``trivial at
  one maximal ideal, hence at all of them by $S_N$'', has no counterpart.
  The route is then: flattening, Theorem~\ref{main} over $D$, and
  push-forward from $D$ to $R$.
\item Correction to an earlier remark: LQ X.5.11 (Milnor's theorem) is
  about Milnor squares, not Zariski squares, so it is not the patching used
  here. Zariski gluing of modules is LQ XV.4, and in Unit it is just gluing
  of families.
\item Inputs to formalise: LQ X.5.5 (via $\lambda_t$ and $\gamma_t$, LQ
  Problem X-2, using exterior powers from \texttt{sag/exterior}); LQ
  XVI.3.1--3.4 (short); Whitehead's identity; Lemmas~\ref{efface}
  and~\ref{split} (elementary); the descent of weight-$0$ families.
\end{enumerate}

\end{document}
